% Generated by thesis/ch7/exp2/signal_tables.py (via make_assets.py) -- do not edit by hand.
\begin{sidewaystable}
  \centering
  \footnotesize
  \setlength{\tabcolsep}{2.6pt}
  \begin{tabularx}{\textheight}{@{}>{\raggedright\arraybackslash}X*{10}{c}>{\raggedright\arraybackslash}p{4.4cm}@{}}
    \toprule
    Signal & s0 & s1 & s2 & s3 & s4 & s5 & s6 & s7 & s8 & s9 & Notes (realisation) \\
    \midrule
    \multicolumn{12}{@{}l}{\textit{Linear core (slopes and one-hot offsets), strongest in the data first}} \\
    \textbf{Marital status and household role}: low income likelier outside marriage: husbands and wives sit well below the unmarried & \checkmark & \checkmark & \checkmark & \checkmark & \checkmark & \checkmark & \checkmark & \checkmark & \checkmark & \checkmark & edge 2, edge+units 8; one-level units in 5 \\
    \textbf{Occupation and work class}: executive, managerial and professional work sit below; service and handling work and missing job records above & \checkmark & \checkmark & \checkmark & \checkmark & \checkmark & \checkmark & \checkmark & \checkmark & \checkmark & \checkmark & edge+units 10; one-level units in 9 \\
    \textbf{Education (years)}: each extra year of schooling lowers the chance of a low income & \checkmark & \checkmark & \checkmark & \checkmark & \checkmark & \checkmark & \checkmark & \checkmark & \checkmark & \checkmark & edge 9, edge+units 1 \\
    \textbf{Capital gain (slope)}: capital gains lower the chance of a low income & \checkmark & \checkmark & \checkmark & \checkmark & \checkmark & \checkmark & \checkmark & \checkmark & \checkmark & \checkmark & edge+units 9, units only 1 \\
    \textbf{Age (slope)}: older people are less likely to have a low income & \checkmark & \checkmark & \checkmark & \checkmark & \checkmark & \checkmark & \checkmark & \checkmark & \checkmark & \checkmark & edge 4, edge+units 5, units only 1 \\
    \textbf{Hours per week (slope)}: longer working weeks lower the chance of a low income & \checkmark & \checkmark & \checkmark & \checkmark & \checkmark & \checkmark & \checkmark & \checkmark & \checkmark & \checkmark & edge 8, edge+units 1, units only 1 \\
    \textbf{Capital loss (slope)}: recorded capital losses lower the chance of a low income & \checkmark & \checkmark & \checkmark & \checkmark & \checkmark & \checkmark & \checkmark & \checkmark & \checkmark & \checkmark & edge 10 \\
    \textbf{Sex}: women are more likely to have a low income than otherwise similar men & \checkmark & \checkmark & \checkmark & \checkmark & \checkmark & \checkmark & \checkmark & \checkmark & \checkmark & \checkmark & edge 9, units only 1 \\
    \textbf{Race}: small offsets by recorded race & w & w & \checkmark & w &  & w & w & w & w & w & edge 1 \\
    \textbf{Education levels}: offsets for single education levels beyond the years-of-schooling line & w & \checkmark & \checkmark & \checkmark & \checkmark & \checkmark & \checkmark & \checkmark & \checkmark & \checkmark & edge 3, edge+units 6; one-level units in 3 \\
    \textbf{Country of birth}: offsets for single countries of birth, mostly small groups &  & \checkmark & \checkmark & \checkmark & \checkmark & \checkmark & w & \checkmark & w & \checkmark & edge+units 7; one-level units in 3; small-group offsets only in 3 \\
    \midrule
    \multicolumn{12}{@{}l}{\textit{Non-linear: bends and interactions, strongest in the data first}} \\
    \textbf{Capital-gain threshold}: past about \$6,000 of capital gain the chance of a low income drops steeply & w & S2 & S2 & S2 &  & S3 & H & G & S3 & S2 & S2 4, S3 2, H 1, G 1; 1-2 parts \\
    \textbf{Age bend}: the fall with age is steep for the young and turns back later in life &  &  &  & H & G & S2 & H & G & G &  & G 3, H 2, S2 1; 1 part \\
    \textbf{Hours bend}: short weeks raise the chance of a low income more than long weeks lower it &  &  &  &  &  & H &  &  &  & H & H 2; 1 part \\
    \midrule
    \multicolumn{12}{@{}l}{\textit{Relevant in the data but in no network}} \\
    \textbf{Capital-loss bend}: bends in the capital-loss effect across loss amounts &  &  &  &  &  &  &  &  &  &  &  \\
    \textbf{Census weight (slope)}: the survey sampling weight, not a personal attribute & w & w & w & w & w & w & w & w &  & w & extreme rows only in 1 \\
    \textbf{Age $\times$ marriage}: the age effect is weaker inside marriage &  &  &  &  &  &  &  &  &  &  &  \\
    \textbf{Census weight bend}: bends in the survey sampling weight &  &  &  &  &  &  &  &  &  &  &  \\
    \textbf{Age $\times$ hours}: the hours effect is weaker for older people &  &  &  &  &  &  &  &  &  &  &  \\
    \textbf{Capital gain $\times$ marriage}: the capital-gain threshold acts more strongly for the married &  &  &  &  &  &  &  &  &  &  &  \\
    \textbf{Sex $\times$ marriage}: the gap between women and men is larger outside marriage &  &  &  &  &  &  &  &  &  &  &  \\
    \bottomrule
  \end{tabularx}
  \caption{Adult: signals (rows) by seed (columns). The positive class is an income of at most \$50K, so a signal that raises the score makes a low income likelier. Each network's logit is projected on a basis of data functions (\texttt{experiments/adult\_signals}): linear terms per variable or one-hot group, bends for continuous inputs, and interactions the explanations' units read together; the projection explains 0.998 to 1.000 of the logit's variance (training rows). A single one-hot switching unit counts as linear. A seed \emph{uses} a signal when its component, clipped to its 1st to 99th percentiles, has a standard deviation of at least 0.10 logit. \checkmark: used; H, one hinge; G, one gate; S$k$, split across $k$ units; w: weak (0.03 to 0.10 logit). Notes: routes for linear signals (edge, units) and realisation counts with explanation parts. Rows are ordered by the data's own strength (cross-validated log-loss gain in a ridge fit on the same basis).}
  \label{tab:exp2-signals-adult}
\end{sidewaystable}
